\section{Discussion}
\label{sec:discussion}

\paragraph{Two sources of character decodability.}
The data separate two sources of linearly decodable character information. One is stored with the
word: present at its tokens in a sentence that asks nothing, largely already in the token embedding,
and roughly constant across a more than 20-fold range of model size. The other appears at the answer position
only when the property is asked for, mostly in response to that specific property, and in the late
half of the network. Only the second grows with scale, so the common observation that larger models
``know'' more about spelling at the answer position is, in this battery, an observation about the
second source.

\paragraph{The late signal is not a relay of the store.}
If the question merely relayed what a linear probe reads at the word's tokens, the answer position
should not systematically exceed them. It falls short of them in small models and exceeds them in
large ones, and the difference grows with scale for both tasks where it can be measured. The excess may
come from combining information spread over the word's several tokens, from computation the
question triggers, or from re-encoding the property into a form a linear probe reads more easily. A
linear probe cannot tell these apart, and we avoid the word ``routing'' because it implies the first
reading only.

\paragraph{The early phase.}
The early bump in the question arm is similar at every scale and is consistent with the word's
spelling being carried to the adjacent answer token by early attention. It explains why some
question-free signal survives at zero shots and why small models' held-out layer selection can land
in the first few layers. It is also a reason to report layer-restricted analyses: a peak over all
layers mixes a scale-invariant early signal with a scale-dependent late one.

\paragraph{What the controls changed.}
Retained question-free signal measured above the majority class was largely nuisance; count tasks
that appeared to show a computation gap failed a word-only floor, and length matching alone did not
fix them; a strong raw-scale narrowing of the \texttt{last\_letter} readout gap disappeared on the
logit scale. A question placed before the word raised the word tokens' \texttt{last\_letter}
decodability in all 16 models in our main runs, but the pre-registered follow-up shows that this is
prefix sensitivity, not anticipation: changing \texttt{The word "} to \texttt{In the word "} moves
the same tokens by a similar amount, and a question adds nothing measurable beyond that. None of
these controls is expensive, and each of them changed a conclusion.
